% Supplementary material for GADAR.
%
% Compiles two ways:
%   (a) standalone, producing a separate supplementary PDF: pdflatex appendix
%   (b) % Supplementary material for GADAR.
%
% Compiles two ways:
%   (a) standalone, producing a separate supplementary PDF: pdflatex appendix
%   (b) % Supplementary material for GADAR.
%
% Compiles two ways:
%   (a) standalone, producing a separate supplementary PDF: pdflatex appendix
%   (b) % Supplementary material for GADAR.
%
% Compiles two ways:
%   (a) standalone, producing a separate supplementary PDF: pdflatex appendix
%   (b) \input{appendix.tex} just before \end{document} of gadar.tex
% The guard below detects which, so the same file serves both.
%
% Everything here is promised by an answer in the reproducibility checklist;
% the mapping is in Table~\ref{tab:checklist-map}.

\makeatletter
\@ifundefined{gadar}{
  \newif\ifappendixStandalone
  \appendixStandalonetrue
}{
  \newif\ifappendixStandalone
  \appendixStandalonefalse
}
\makeatother

\ifappendixStandalone
\documentclass[letterpaper]{article}
\usepackage{aaai2027}
\usepackage[hyphens]{url}
\usepackage{graphicx}
\urlstyle{rm}
\def\UrlFont{\rm}
\usepackage{natbib}
\usepackage{caption}
\usepackage{booktabs}
\usepackage{multirow}
\usepackage{amsfonts}
\usepackage{amsmath}
\usepackage{algorithm}
\usepackage{algorithmic}
\frenchspacing
\setlength{\pdfpagewidth}{8.5in}
\setlength{\pdfpageheight}{11in}
\setcounter{secnumdepth}{0}

\newcommand{\gadar}{GADAR}
\newcommand{\gadarbind}{\gadar{}-BIND}
\newcommand{\union}{UNION}

\title{GADAR: Supplementary Material}
\author{Anonymous Submission}
\affiliations{Anonymized for review}

\begin{document}
\maketitle
\fi

\appendix
\section{Supplementary Material}

This appendix supplies the detail promised by our reproducibility checklist
answers. Table~\ref{tab:checklist-map} maps each answer to the section that
discharges it.

\begin{table}[h]
\small
\centering
\caption{Checklist answers and where they are discharged.}
\label{tab:checklist-map}
\begin{tabular}{@{}llp{2.9cm}@{}}
\toprule
Item & Answer & Section \\
\midrule
1.1 outline / pseudocode        & yes     & \ref{app:pseudocode} \\
4.2 hyperparameter search       & partial & \ref{app:hyper} \\
4.5 code released on publication& yes     & \ref{app:code} \\
4.6 implementation comments     & partial & \ref{app:code} \\
4.7 seeding                     & yes     & \ref{app:seeds} \\
4.8 computing infrastructure    & partial & \ref{app:compute} \\
4.9 evaluation metrics          & yes     & \ref{app:metrics} \\
4.10 number of runs             & yes     & \ref{app:seeds} \\
4.13 final hyperparameters      & yes     & \ref{app:hyper} \\
\bottomrule
\end{tabular}
\end{table}

\subsection{Code and Data Availability}
\label{app:code}

The implementation is a single Python package built on PyTorch and PyTorch
Geometric, using PDDLGym to parse domains and problems and Fast Downward
\citep{fast_downward} to generate training plans. The full source, the
experiment configuration files that reproduce every number in the paper, and
the analysis scripts will be released publicly under a permissive license on
publication. They are not attached to this submission.

Each experiment is a YAML configuration file naming the training domains, the
featurization rung, and the evaluation domains; a run is a single invocation
of the entry point with one such file, so the configuration files are a
complete record of what was run. Files implementing the components introduced
here, in particular graph construction, the lifted domain layer, the binding
layer, and the chain layer, carry comments naming the paper section each step
corresponds to. Files inherited unchanged from the per-domain policy we build
on \citep{gabar} are not annotated in that way.

\subsection{Domains and Data}
\label{app:data}

We use the eight benchmark domains of \citet{gabar}, generated with
\citet{pddl-generators}. Training data is the set of states along a
satisficing plan for each small instance, labelled with the action the planner
took. Table~\ref{tab:data} gives the sizes. Instances the planner fails to
solve within its limit are skipped, so the collected count can be below the
number of instances the generator produced.

\begin{table}[h]
\small
\centering
\caption{Data per domain. Problems and states are training instances and the
labelled states collected from their plans; test problems are the domain's
full test set (our hard split).}
\label{tab:data}
\begin{tabular}{@{}lrrr@{}}
\toprule
Domain & Problems & States & Test problems \\
\midrule
Blocks     & 200 & 4{,}216 & 200 \\
Gripper    & 147 & 7{,}287 & 173 \\
Miconic    & 228 & 5{,}453 & 119 \\
Visitall   & 125 & 3{,}654 & 50 \\
Grid       & 192 & 6{,}880 & 48 \\
Logistics  & 156 & 6{,}732 & 96 \\
Spanner    & 234 & 3{,}801 & 96 \\
Rovers     & 312 & 3{,}001 & 54 \\
\midrule
Total      & 1{,}594 & 41{,}024 & 836 \\
\bottomrule
\end{tabular}
\end{table}

Ten percent of the states of each domain are held out as a validation split,
stratified by domain so every training domain is represented. The split is
deterministic given the seed. In the leave-one-domain-out experiments the
held-out domain contributes no training or validation data at all: its plans
are never collected for that run.

\subsection{Graph Construction}
\label{app:pseudocode}

Algorithm~\ref{alg:graph} states the construction of
$\mathcal{G} = (V, E, X, R)$ for one task, in the notation of the
representation section. Lines marked $\dagger$ are computed once per domain
and reused for every state of that domain; the rest are per state. The three
rungs of the ladder are prefixes of this procedure: \union{} runs lines
1 to 3 with symbol-identity features in place of $X$, \gadarbind{} runs lines
1 to 7, and \gadar{} runs all of it.

\begin{algorithm}[h]
\caption{Construction of the \gadar{} graph for $(s, G, D)$}
\label{alg:graph}
\begin{algorithmic}[1]
\STATE $V \leftarrow O \cup L \cup A \cup \{g\}$ where $L$ holds the ground
  literals of $s$ and of $G$, the latter flagged as goal atoms
\STATE add $E_{\mathrm{lit}}$: each literal to its argument objects, argument
  position on the edge
\STATE add $E_{\mathrm{act}}$: each schema to the objects of its
  \emph{applicable} groundings, parameter slot on the edge
\STATE $\dagger$ compile each declared PDDL type of $D$ into a unary static
  predicate, so typed and predicate-typed domains present one structure
\STATE $\dagger$ $V \leftarrow V \cup \Sigma \cup Q$: a node per predicate
  symbol of $D$, and an occurrence node $q_{a,k}$ per literal occurrence in
  the preconditions and effects of each schema $a$
\STATE $\dagger$ add $E_{\mathrm{occ}}$: each occurrence to its predicate
  symbol and its schema, carrying role $r(q)$ and argument-to-slot pairs
\STATE add $E_{\mathrm{inst}}$ and $E_{\mathrm{bind}}$: each literal to its
  predicate symbol, each object to its type symbol, and $(o, q_{a,k})$
  whenever an applicable grounding of $a$ puts $o$ at a slot that occurrence
  $k$ constrains
\STATE $\dagger$ add $E_{\mathrm{chain}}$: \textit{enables} between an add and
  a pre occurrence of a shared predicate, \textit{threatens} likewise for del
\STATE add $E_{\mathrm{eff}}$: from each applicable grounding, link its
  objects to any unsatisfied goal atom it adds and any true atom it deletes
\STATE compute $\mathrm{gdist}(a)$ by breadth-first search over
  \textit{enables} from the schemas adding a goal predicate; mark objects
  occurring in an open goal atom
\STATE emit $X$ and $R$ as indicator vectors; widths depend only on the global
  arity caps $(k_p, k_a)$, never on $|A|$, $|P|$, $|T|$ or $|O|$
\end{algorithmic}
\end{algorithm}

Decoding is unchanged from the per-domain policy: a GRU scores the schema
nodes against its hidden state, takes the best, then scores the object nodes
once per argument of the selected schema, feeding each choice back as the next
input. Search over these choices is a beam of width 2 over schemas and 3 over
objects per slot.

\subsection{Hyperparameters}
\label{app:hyper}

Table~\ref{tab:hyper} lists every setting used. Encoder and decoder
hyperparameters are inherited from \citet{gabar} without modification, which
is deliberate: the paper's claims are about the representation, and re-tuning
the network per rung would confound them. We therefore ran no hyperparameter
search over the architecture, and item 4.2 of the checklist is answered
``partial'' for that reason. The two settings we did choose are the global
arity caps $(k_p, k_a)$, set to the maxima over the training domains so that a
held-out domain with wider symbols still featurizes at full width, and the
training length, discussed below.

\begin{table}[h]
\small
\centering
\caption{Final hyperparameters. Identical across all three rungs except where
noted.}
\label{tab:hyper}
\begin{tabular}{@{}lr@{}}
\toprule
Setting & Value \\
\midrule
\multicolumn{2}{@{}l}{\emph{Encoder}} \\
Message-passing rounds $L$        & 9 \\
Representation size               & 64 \\
Attention heads                   & 1 \\
MLP layers per update block       & 2 \\
Global node                       & enabled \\
\midrule
\multicolumn{2}{@{}l}{\emph{Decoder}} \\
GRU layers                        & 3 \\
Beam width, schemas               & 2 \\
Beam width, objects per slot      & 3 \\
\midrule
\multicolumn{2}{@{}l}{\emph{Optimization}} \\
Optimizer                         & Adam \\
Learning rate                     & $5 \times 10^{-4}$ \\
Batch size                        & 64 \\
Dropout                           & 0.0 \\
Attention dropout                 & 0.0 \\
Weight decay                      & 0.0 \\
Epochs, structural rungs          & 500 \\
Epochs, \union{}                  & 100 \\
\midrule
\multicolumn{2}{@{}l}{\emph{Evaluation}} \\
Execution bound (steps)           & 500 \\
Search at test time               & none (greedy) \\
Revisit monitor                   & enabled \\
Training plan generator           & \texttt{fd-lama-first} \\
Plan-quality reference            & \texttt{fd-lama-first} \\
\bottomrule
\end{tabular}
\end{table}

\subsection{Seeds and Number of Runs}
\label{app:seeds}

A run is seeded once, and that seed initializes the network, the training
shuffle, the train and validation split, and the revisit monitor's fallback
sampling. The seed is recorded alongside every checkpoint, and the checkpoint
directory is keyed by the training domain set, the seed, the featurization,
and the resulting input widths, so runs that differ in any of these cannot
overwrite one another.

Each reported configuration is run for three seeds and the results averaged.
We report means without dispersion: three runs are too few for a meaningful
interval, and the differences the paper turns on, for instance 4\% against
68\% zero-shot coverage on the hard split of held-out Visitall, are far larger
than the spread we observe. No statistical test is applied, which is item 4.12
of the checklist.

\subsection{Evaluation Metrics}
\label{app:metrics}

Let $\mathcal{T}$ be a test set and $\mathrm{solved}(\pi) \subseteq
\mathcal{T}$ the instances on which policy $\pi$ reaches a goal state within
the execution bound. \emph{Coverage} is
\begin{equation*}
C(\pi) \;=\; \frac{|\mathrm{solved}(\pi)|}{|\mathcal{T}|} \times 100 .
\end{equation*}
Execution is greedy and search-free: at each state the highest-ranked
applicable action the decoder constructs is executed, and a rollout ends at
the goal, at a dead end, or at the bound. A rollout that proposes an
inapplicable action falls back to the next candidate in the beam; if none is
applicable the rollout fails.

\emph{Plan quality} compares plan lengths on the instances both the policy and
the reference planner solve. With $|\pi(t)|$ the policy's plan length on
instance $t$ and $|\rho(t)|$ the reference planner's,
\begin{equation*}
P(\pi) \;=\; \frac{\sum_{t \in \mathrm{solved}(\pi) \cap \mathrm{solved}(\rho)} |\rho(t)|}
                 {\sum_{t \in \mathrm{solved}(\pi) \cap \mathrm{solved}(\rho)} |\pi(t)|} .
\end{equation*}
Higher is better and $P > 1$ means plans shorter than the reference. $P$ is
conditioned on what a system solves, so a system with higher coverage is
scored on harder instances and $P$ must be read within a column rather than
across columns. Coverage is the primary metric for that reason.

Coverage is measured with a revisit monitor, which resamples among the
decoder's applicable proposals when a greedy rollout re-enters a state it has
already visited. The monitor is part of the execution harness and is applied
identically to every system, including the untrained control, which is what
makes the control informative: it measures what the harness achieves when the
learned component contributes nothing.

\subsection{Compute Environment}
\label{app:compute}

All training and all evaluation ran on a single NVIDIA A40 GPU per run, one
run per GPU. Training a single model on all eight domains takes seven to eight
hours; for comparison, the per-domain models of \citet{gabar} take roughly two
hours each, so eight of them cost about sixteen hours and produce eight
artifacts rather than one.

The implementation uses PyTorch and PyTorch Geometric on Linux; a pinned
dependency specification ships with the code release.

\subsection{Renaming Invariance}
\label{app:invariance}

The claim that \gadar{} is invariant to renaming is enforced by construction
rather than learned, and it is checked by a unit test that runs on every
change to the representation: given a domain and a bijective renaming of its
predicates, action schemas, types, and objects, the test builds the graph for
the original and the renamed task and asserts that the node feature matrix,
the edge feature matrix, and the edge lists are identical. The test covers
both typing conventions in the benchmark suite, domains that declare PDDL
types and domains that spell types as ordinary unary predicates, and asserts
that the two produce the same structure after type compilation.

\ifappendixStandalone
\bibliography{aaai2027}
\end{document}
\fi
 just before \end{document} of gadar.tex
% The guard below detects which, so the same file serves both.
%
% Everything here is promised by an answer in the reproducibility checklist;
% the mapping is in Table~\ref{tab:checklist-map}.

\makeatletter
\@ifundefined{gadar}{
  \newif\ifappendixStandalone
  \appendixStandalonetrue
}{
  \newif\ifappendixStandalone
  \appendixStandalonefalse
}
\makeatother

\ifappendixStandalone
\documentclass[letterpaper]{article}
\usepackage{aaai2027}
\usepackage[hyphens]{url}
\usepackage{graphicx}
\urlstyle{rm}
\def\UrlFont{\rm}
\usepackage{natbib}
\usepackage{caption}
\usepackage{booktabs}
\usepackage{multirow}
\usepackage{amsfonts}
\usepackage{amsmath}
\usepackage{algorithm}
\usepackage{algorithmic}
\frenchspacing
\setlength{\pdfpagewidth}{8.5in}
\setlength{\pdfpageheight}{11in}
\setcounter{secnumdepth}{0}

\newcommand{\gadar}{GADAR}
\newcommand{\gadarbind}{\gadar{}-BIND}
\newcommand{\union}{UNION}

\title{GADAR: Supplementary Material}
\author{Anonymous Submission}
\affiliations{Anonymized for review}

\begin{document}
\maketitle
\fi

\appendix
\section{Supplementary Material}

This appendix supplies the detail promised by our reproducibility checklist
answers. Table~\ref{tab:checklist-map} maps each answer to the section that
discharges it.

\begin{table}[h]
\small
\centering
\caption{Checklist answers and where they are discharged.}
\label{tab:checklist-map}
\begin{tabular}{@{}llp{2.9cm}@{}}
\toprule
Item & Answer & Section \\
\midrule
1.1 outline / pseudocode        & yes     & \ref{app:pseudocode} \\
4.2 hyperparameter search       & partial & \ref{app:hyper} \\
4.5 code released on publication& yes     & \ref{app:code} \\
4.6 implementation comments     & partial & \ref{app:code} \\
4.7 seeding                     & yes     & \ref{app:seeds} \\
4.8 computing infrastructure    & partial & \ref{app:compute} \\
4.9 evaluation metrics          & yes     & \ref{app:metrics} \\
4.10 number of runs             & yes     & \ref{app:seeds} \\
4.13 final hyperparameters      & yes     & \ref{app:hyper} \\
\bottomrule
\end{tabular}
\end{table}

\subsection{Code and Data Availability}
\label{app:code}

The implementation is a single Python package built on PyTorch and PyTorch
Geometric, using PDDLGym to parse domains and problems and Fast Downward
\citep{fast_downward} to generate training plans. The full source, the
experiment configuration files that reproduce every number in the paper, and
the analysis scripts will be released publicly under a permissive license on
publication. They are not attached to this submission.

Each experiment is a YAML configuration file naming the training domains, the
featurization rung, and the evaluation domains; a run is a single invocation
of the entry point with one such file, so the configuration files are a
complete record of what was run. Files implementing the components introduced
here, in particular graph construction, the lifted domain layer, the binding
layer, and the chain layer, carry comments naming the paper section each step
corresponds to. Files inherited unchanged from the per-domain policy we build
on \citep{gabar} are not annotated in that way.

\subsection{Domains and Data}
\label{app:data}

We use the eight benchmark domains of \citet{gabar}, generated with
\citet{pddl-generators}. Training data is the set of states along a
satisficing plan for each small instance, labelled with the action the planner
took. Table~\ref{tab:data} gives the sizes. Instances the planner fails to
solve within its limit are skipped, so the collected count can be below the
number of instances the generator produced.

\begin{table}[h]
\small
\centering
\caption{Data per domain. Problems and states are training instances and the
labelled states collected from their plans; test problems are the domain's
full test set (our hard split).}
\label{tab:data}
\begin{tabular}{@{}lrrr@{}}
\toprule
Domain & Problems & States & Test problems \\
\midrule
Blocks     & 200 & 4{,}216 & 200 \\
Gripper    & 147 & 7{,}287 & 173 \\
Miconic    & 228 & 5{,}453 & 119 \\
Visitall   & 125 & 3{,}654 & 50 \\
Grid       & 192 & 6{,}880 & 48 \\
Logistics  & 156 & 6{,}732 & 96 \\
Spanner    & 234 & 3{,}801 & 96 \\
Rovers     & 312 & 3{,}001 & 54 \\
\midrule
Total      & 1{,}594 & 41{,}024 & 836 \\
\bottomrule
\end{tabular}
\end{table}

Ten percent of the states of each domain are held out as a validation split,
stratified by domain so every training domain is represented. The split is
deterministic given the seed. In the leave-one-domain-out experiments the
held-out domain contributes no training or validation data at all: its plans
are never collected for that run.

\subsection{Graph Construction}
\label{app:pseudocode}

Algorithm~\ref{alg:graph} states the construction of
$\mathcal{G} = (V, E, X, R)$ for one task, in the notation of the
representation section. Lines marked $\dagger$ are computed once per domain
and reused for every state of that domain; the rest are per state. The three
rungs of the ladder are prefixes of this procedure: \union{} runs lines
1 to 3 with symbol-identity features in place of $X$, \gadarbind{} runs lines
1 to 7, and \gadar{} runs all of it.

\begin{algorithm}[h]
\caption{Construction of the \gadar{} graph for $(s, G, D)$}
\label{alg:graph}
\begin{algorithmic}[1]
\STATE $V \leftarrow O \cup L \cup A \cup \{g\}$ where $L$ holds the ground
  literals of $s$ and of $G$, the latter flagged as goal atoms
\STATE add $E_{\mathrm{lit}}$: each literal to its argument objects, argument
  position on the edge
\STATE add $E_{\mathrm{act}}$: each schema to the objects of its
  \emph{applicable} groundings, parameter slot on the edge
\STATE $\dagger$ compile each declared PDDL type of $D$ into a unary static
  predicate, so typed and predicate-typed domains present one structure
\STATE $\dagger$ $V \leftarrow V \cup \Sigma \cup Q$: a node per predicate
  symbol of $D$, and an occurrence node $q_{a,k}$ per literal occurrence in
  the preconditions and effects of each schema $a$
\STATE $\dagger$ add $E_{\mathrm{occ}}$: each occurrence to its predicate
  symbol and its schema, carrying role $r(q)$ and argument-to-slot pairs
\STATE add $E_{\mathrm{inst}}$ and $E_{\mathrm{bind}}$: each literal to its
  predicate symbol, each object to its type symbol, and $(o, q_{a,k})$
  whenever an applicable grounding of $a$ puts $o$ at a slot that occurrence
  $k$ constrains
\STATE $\dagger$ add $E_{\mathrm{chain}}$: \textit{enables} between an add and
  a pre occurrence of a shared predicate, \textit{threatens} likewise for del
\STATE add $E_{\mathrm{eff}}$: from each applicable grounding, link its
  objects to any unsatisfied goal atom it adds and any true atom it deletes
\STATE compute $\mathrm{gdist}(a)$ by breadth-first search over
  \textit{enables} from the schemas adding a goal predicate; mark objects
  occurring in an open goal atom
\STATE emit $X$ and $R$ as indicator vectors; widths depend only on the global
  arity caps $(k_p, k_a)$, never on $|A|$, $|P|$, $|T|$ or $|O|$
\end{algorithmic}
\end{algorithm}

Decoding is unchanged from the per-domain policy: a GRU scores the schema
nodes against its hidden state, takes the best, then scores the object nodes
once per argument of the selected schema, feeding each choice back as the next
input. Search over these choices is a beam of width 2 over schemas and 3 over
objects per slot.

\subsection{Hyperparameters}
\label{app:hyper}

Table~\ref{tab:hyper} lists every setting used. Encoder and decoder
hyperparameters are inherited from \citet{gabar} without modification, which
is deliberate: the paper's claims are about the representation, and re-tuning
the network per rung would confound them. We therefore ran no hyperparameter
search over the architecture, and item 4.2 of the checklist is answered
``partial'' for that reason. The two settings we did choose are the global
arity caps $(k_p, k_a)$, set to the maxima over the training domains so that a
held-out domain with wider symbols still featurizes at full width, and the
training length, discussed below.

\begin{table}[h]
\small
\centering
\caption{Final hyperparameters. Identical across all three rungs except where
noted.}
\label{tab:hyper}
\begin{tabular}{@{}lr@{}}
\toprule
Setting & Value \\
\midrule
\multicolumn{2}{@{}l}{\emph{Encoder}} \\
Message-passing rounds $L$        & 9 \\
Representation size               & 64 \\
Attention heads                   & 1 \\
MLP layers per update block       & 2 \\
Global node                       & enabled \\
\midrule
\multicolumn{2}{@{}l}{\emph{Decoder}} \\
GRU layers                        & 3 \\
Beam width, schemas               & 2 \\
Beam width, objects per slot      & 3 \\
\midrule
\multicolumn{2}{@{}l}{\emph{Optimization}} \\
Optimizer                         & Adam \\
Learning rate                     & $5 \times 10^{-4}$ \\
Batch size                        & 64 \\
Dropout                           & 0.0 \\
Attention dropout                 & 0.0 \\
Weight decay                      & 0.0 \\
Epochs, structural rungs          & 500 \\
Epochs, \union{}                  & 100 \\
\midrule
\multicolumn{2}{@{}l}{\emph{Evaluation}} \\
Execution bound (steps)           & 500 \\
Search at test time               & none (greedy) \\
Revisit monitor                   & enabled \\
Training plan generator           & \texttt{fd-lama-first} \\
Plan-quality reference            & \texttt{fd-lama-first} \\
\bottomrule
\end{tabular}
\end{table}

\subsection{Seeds and Number of Runs}
\label{app:seeds}

A run is seeded once, and that seed initializes the network, the training
shuffle, the train and validation split, and the revisit monitor's fallback
sampling. The seed is recorded alongside every checkpoint, and the checkpoint
directory is keyed by the training domain set, the seed, the featurization,
and the resulting input widths, so runs that differ in any of these cannot
overwrite one another.

Each reported configuration is run for three seeds and the results averaged.
We report means without dispersion: three runs are too few for a meaningful
interval, and the differences the paper turns on, for instance 4\% against
68\% zero-shot coverage on the hard split of held-out Visitall, are far larger
than the spread we observe. No statistical test is applied, which is item 4.12
of the checklist.

\subsection{Evaluation Metrics}
\label{app:metrics}

Let $\mathcal{T}$ be a test set and $\mathrm{solved}(\pi) \subseteq
\mathcal{T}$ the instances on which policy $\pi$ reaches a goal state within
the execution bound. \emph{Coverage} is
\begin{equation*}
C(\pi) \;=\; \frac{|\mathrm{solved}(\pi)|}{|\mathcal{T}|} \times 100 .
\end{equation*}
Execution is greedy and search-free: at each state the highest-ranked
applicable action the decoder constructs is executed, and a rollout ends at
the goal, at a dead end, or at the bound. A rollout that proposes an
inapplicable action falls back to the next candidate in the beam; if none is
applicable the rollout fails.

\emph{Plan quality} compares plan lengths on the instances both the policy and
the reference planner solve. With $|\pi(t)|$ the policy's plan length on
instance $t$ and $|\rho(t)|$ the reference planner's,
\begin{equation*}
P(\pi) \;=\; \frac{\sum_{t \in \mathrm{solved}(\pi) \cap \mathrm{solved}(\rho)} |\rho(t)|}
                 {\sum_{t \in \mathrm{solved}(\pi) \cap \mathrm{solved}(\rho)} |\pi(t)|} .
\end{equation*}
Higher is better and $P > 1$ means plans shorter than the reference. $P$ is
conditioned on what a system solves, so a system with higher coverage is
scored on harder instances and $P$ must be read within a column rather than
across columns. Coverage is the primary metric for that reason.

Coverage is measured with a revisit monitor, which resamples among the
decoder's applicable proposals when a greedy rollout re-enters a state it has
already visited. The monitor is part of the execution harness and is applied
identically to every system, including the untrained control, which is what
makes the control informative: it measures what the harness achieves when the
learned component contributes nothing.

\subsection{Compute Environment}
\label{app:compute}

All training and all evaluation ran on a single NVIDIA A40 GPU per run, one
run per GPU. Training a single model on all eight domains takes seven to eight
hours; for comparison, the per-domain models of \citet{gabar} take roughly two
hours each, so eight of them cost about sixteen hours and produce eight
artifacts rather than one.

The implementation uses PyTorch and PyTorch Geometric on Linux; a pinned
dependency specification ships with the code release.

\subsection{Renaming Invariance}
\label{app:invariance}

The claim that \gadar{} is invariant to renaming is enforced by construction
rather than learned, and it is checked by a unit test that runs on every
change to the representation: given a domain and a bijective renaming of its
predicates, action schemas, types, and objects, the test builds the graph for
the original and the renamed task and asserts that the node feature matrix,
the edge feature matrix, and the edge lists are identical. The test covers
both typing conventions in the benchmark suite, domains that declare PDDL
types and domains that spell types as ordinary unary predicates, and asserts
that the two produce the same structure after type compilation.

\ifappendixStandalone
\bibliography{aaai2027}
\end{document}
\fi
 just before \end{document} of gadar.tex
% The guard below detects which, so the same file serves both.
%
% Everything here is promised by an answer in the reproducibility checklist;
% the mapping is in Table~\ref{tab:checklist-map}.

\makeatletter
\@ifundefined{gadar}{
  \newif\ifappendixStandalone
  \appendixStandalonetrue
}{
  \newif\ifappendixStandalone
  \appendixStandalonefalse
}
\makeatother

\ifappendixStandalone
\documentclass[letterpaper]{article}
\usepackage{aaai2027}
\usepackage[hyphens]{url}
\usepackage{graphicx}
\urlstyle{rm}
\def\UrlFont{\rm}
\usepackage{natbib}
\usepackage{caption}
\usepackage{booktabs}
\usepackage{multirow}
\usepackage{amsfonts}
\usepackage{amsmath}
\usepackage{algorithm}
\usepackage{algorithmic}
\frenchspacing
\setlength{\pdfpagewidth}{8.5in}
\setlength{\pdfpageheight}{11in}
\setcounter{secnumdepth}{0}

\newcommand{\gadar}{GADAR}
\newcommand{\gadarbind}{\gadar{}-BIND}
\newcommand{\union}{UNION}

\title{GADAR: Supplementary Material}
\author{Anonymous Submission}
\affiliations{Anonymized for review}

\begin{document}
\maketitle
\fi

\appendix
\section{Supplementary Material}

This appendix supplies the detail promised by our reproducibility checklist
answers. Table~\ref{tab:checklist-map} maps each answer to the section that
discharges it.

\begin{table}[h]
\small
\centering
\caption{Checklist answers and where they are discharged.}
\label{tab:checklist-map}
\begin{tabular}{@{}llp{2.9cm}@{}}
\toprule
Item & Answer & Section \\
\midrule
1.1 outline / pseudocode        & yes     & \ref{app:pseudocode} \\
4.2 hyperparameter search       & partial & \ref{app:hyper} \\
4.5 code released on publication& yes     & \ref{app:code} \\
4.6 implementation comments     & partial & \ref{app:code} \\
4.7 seeding                     & yes     & \ref{app:seeds} \\
4.8 computing infrastructure    & partial & \ref{app:compute} \\
4.9 evaluation metrics          & yes     & \ref{app:metrics} \\
4.10 number of runs             & yes     & \ref{app:seeds} \\
4.13 final hyperparameters      & yes     & \ref{app:hyper} \\
\bottomrule
\end{tabular}
\end{table}

\subsection{Code and Data Availability}
\label{app:code}

The implementation is a single Python package built on PyTorch and PyTorch
Geometric, using PDDLGym to parse domains and problems and Fast Downward
\citep{fast_downward} to generate training plans. The full source, the
experiment configuration files that reproduce every number in the paper, and
the analysis scripts will be released publicly under a permissive license on
publication. They are not attached to this submission.

Each experiment is a YAML configuration file naming the training domains, the
featurization rung, and the evaluation domains; a run is a single invocation
of the entry point with one such file, so the configuration files are a
complete record of what was run. Files implementing the components introduced
here, in particular graph construction, the lifted domain layer, the binding
layer, and the chain layer, carry comments naming the paper section each step
corresponds to. Files inherited unchanged from the per-domain policy we build
on \citep{gabar} are not annotated in that way.

\subsection{Domains and Data}
\label{app:data}

We use the eight benchmark domains of \citet{gabar}, generated with
\citet{pddl-generators}. Training data is the set of states along a
satisficing plan for each small instance, labelled with the action the planner
took. Table~\ref{tab:data} gives the sizes. Instances the planner fails to
solve within its limit are skipped, so the collected count can be below the
number of instances the generator produced.

\begin{table}[h]
\small
\centering
\caption{Data per domain. Problems and states are training instances and the
labelled states collected from their plans; test problems are the domain's
full test set (our hard split).}
\label{tab:data}
\begin{tabular}{@{}lrrr@{}}
\toprule
Domain & Problems & States & Test problems \\
\midrule
Blocks     & 200 & 4{,}216 & 200 \\
Gripper    & 147 & 7{,}287 & 173 \\
Miconic    & 228 & 5{,}453 & 119 \\
Visitall   & 125 & 3{,}654 & 50 \\
Grid       & 192 & 6{,}880 & 48 \\
Logistics  & 156 & 6{,}732 & 96 \\
Spanner    & 234 & 3{,}801 & 96 \\
Rovers     & 312 & 3{,}001 & 54 \\
\midrule
Total      & 1{,}594 & 41{,}024 & 836 \\
\bottomrule
\end{tabular}
\end{table}

Ten percent of the states of each domain are held out as a validation split,
stratified by domain so every training domain is represented. The split is
deterministic given the seed. In the leave-one-domain-out experiments the
held-out domain contributes no training or validation data at all: its plans
are never collected for that run.

\subsection{Graph Construction}
\label{app:pseudocode}

Algorithm~\ref{alg:graph} states the construction of
$\mathcal{G} = (V, E, X, R)$ for one task, in the notation of the
representation section. Lines marked $\dagger$ are computed once per domain
and reused for every state of that domain; the rest are per state. The three
rungs of the ladder are prefixes of this procedure: \union{} runs lines
1 to 3 with symbol-identity features in place of $X$, \gadarbind{} runs lines
1 to 7, and \gadar{} runs all of it.

\begin{algorithm}[h]
\caption{Construction of the \gadar{} graph for $(s, G, D)$}
\label{alg:graph}
\begin{algorithmic}[1]
\STATE $V \leftarrow O \cup L \cup A \cup \{g\}$ where $L$ holds the ground
  literals of $s$ and of $G$, the latter flagged as goal atoms
\STATE add $E_{\mathrm{lit}}$: each literal to its argument objects, argument
  position on the edge
\STATE add $E_{\mathrm{act}}$: each schema to the objects of its
  \emph{applicable} groundings, parameter slot on the edge
\STATE $\dagger$ compile each declared PDDL type of $D$ into a unary static
  predicate, so typed and predicate-typed domains present one structure
\STATE $\dagger$ $V \leftarrow V \cup \Sigma \cup Q$: a node per predicate
  symbol of $D$, and an occurrence node $q_{a,k}$ per literal occurrence in
  the preconditions and effects of each schema $a$
\STATE $\dagger$ add $E_{\mathrm{occ}}$: each occurrence to its predicate
  symbol and its schema, carrying role $r(q)$ and argument-to-slot pairs
\STATE add $E_{\mathrm{inst}}$ and $E_{\mathrm{bind}}$: each literal to its
  predicate symbol, each object to its type symbol, and $(o, q_{a,k})$
  whenever an applicable grounding of $a$ puts $o$ at a slot that occurrence
  $k$ constrains
\STATE $\dagger$ add $E_{\mathrm{chain}}$: \textit{enables} between an add and
  a pre occurrence of a shared predicate, \textit{threatens} likewise for del
\STATE add $E_{\mathrm{eff}}$: from each applicable grounding, link its
  objects to any unsatisfied goal atom it adds and any true atom it deletes
\STATE compute $\mathrm{gdist}(a)$ by breadth-first search over
  \textit{enables} from the schemas adding a goal predicate; mark objects
  occurring in an open goal atom
\STATE emit $X$ and $R$ as indicator vectors; widths depend only on the global
  arity caps $(k_p, k_a)$, never on $|A|$, $|P|$, $|T|$ or $|O|$
\end{algorithmic}
\end{algorithm}

Decoding is unchanged from the per-domain policy: a GRU scores the schema
nodes against its hidden state, takes the best, then scores the object nodes
once per argument of the selected schema, feeding each choice back as the next
input. Search over these choices is a beam of width 2 over schemas and 3 over
objects per slot.

\subsection{Hyperparameters}
\label{app:hyper}

Table~\ref{tab:hyper} lists every setting used. Encoder and decoder
hyperparameters are inherited from \citet{gabar} without modification, which
is deliberate: the paper's claims are about the representation, and re-tuning
the network per rung would confound them. We therefore ran no hyperparameter
search over the architecture, and item 4.2 of the checklist is answered
``partial'' for that reason. The two settings we did choose are the global
arity caps $(k_p, k_a)$, set to the maxima over the training domains so that a
held-out domain with wider symbols still featurizes at full width, and the
training length, discussed below.

\begin{table}[h]
\small
\centering
\caption{Final hyperparameters. Identical across all three rungs except where
noted.}
\label{tab:hyper}
\begin{tabular}{@{}lr@{}}
\toprule
Setting & Value \\
\midrule
\multicolumn{2}{@{}l}{\emph{Encoder}} \\
Message-passing rounds $L$        & 9 \\
Representation size               & 64 \\
Attention heads                   & 1 \\
MLP layers per update block       & 2 \\
Global node                       & enabled \\
\midrule
\multicolumn{2}{@{}l}{\emph{Decoder}} \\
GRU layers                        & 3 \\
Beam width, schemas               & 2 \\
Beam width, objects per slot      & 3 \\
\midrule
\multicolumn{2}{@{}l}{\emph{Optimization}} \\
Optimizer                         & Adam \\
Learning rate                     & $5 \times 10^{-4}$ \\
Batch size                        & 64 \\
Dropout                           & 0.0 \\
Attention dropout                 & 0.0 \\
Weight decay                      & 0.0 \\
Epochs, structural rungs          & 500 \\
Epochs, \union{}                  & 100 \\
\midrule
\multicolumn{2}{@{}l}{\emph{Evaluation}} \\
Execution bound (steps)           & 500 \\
Search at test time               & none (greedy) \\
Revisit monitor                   & enabled \\
Training plan generator           & \texttt{fd-lama-first} \\
Plan-quality reference            & \texttt{fd-lama-first} \\
\bottomrule
\end{tabular}
\end{table}

\subsection{Seeds and Number of Runs}
\label{app:seeds}

A run is seeded once, and that seed initializes the network, the training
shuffle, the train and validation split, and the revisit monitor's fallback
sampling. The seed is recorded alongside every checkpoint, and the checkpoint
directory is keyed by the training domain set, the seed, the featurization,
and the resulting input widths, so runs that differ in any of these cannot
overwrite one another.

Each reported configuration is run for three seeds and the results averaged.
We report means without dispersion: three runs are too few for a meaningful
interval, and the differences the paper turns on, for instance 4\% against
68\% zero-shot coverage on the hard split of held-out Visitall, are far larger
than the spread we observe. No statistical test is applied, which is item 4.12
of the checklist.

\subsection{Evaluation Metrics}
\label{app:metrics}

Let $\mathcal{T}$ be a test set and $\mathrm{solved}(\pi) \subseteq
\mathcal{T}$ the instances on which policy $\pi$ reaches a goal state within
the execution bound. \emph{Coverage} is
\begin{equation*}
C(\pi) \;=\; \frac{|\mathrm{solved}(\pi)|}{|\mathcal{T}|} \times 100 .
\end{equation*}
Execution is greedy and search-free: at each state the highest-ranked
applicable action the decoder constructs is executed, and a rollout ends at
the goal, at a dead end, or at the bound. A rollout that proposes an
inapplicable action falls back to the next candidate in the beam; if none is
applicable the rollout fails.

\emph{Plan quality} compares plan lengths on the instances both the policy and
the reference planner solve. With $|\pi(t)|$ the policy's plan length on
instance $t$ and $|\rho(t)|$ the reference planner's,
\begin{equation*}
P(\pi) \;=\; \frac{\sum_{t \in \mathrm{solved}(\pi) \cap \mathrm{solved}(\rho)} |\rho(t)|}
                 {\sum_{t \in \mathrm{solved}(\pi) \cap \mathrm{solved}(\rho)} |\pi(t)|} .
\end{equation*}
Higher is better and $P > 1$ means plans shorter than the reference. $P$ is
conditioned on what a system solves, so a system with higher coverage is
scored on harder instances and $P$ must be read within a column rather than
across columns. Coverage is the primary metric for that reason.

Coverage is measured with a revisit monitor, which resamples among the
decoder's applicable proposals when a greedy rollout re-enters a state it has
already visited. The monitor is part of the execution harness and is applied
identically to every system, including the untrained control, which is what
makes the control informative: it measures what the harness achieves when the
learned component contributes nothing.

\subsection{Compute Environment}
\label{app:compute}

All training and all evaluation ran on a single NVIDIA A40 GPU per run, one
run per GPU. Training a single model on all eight domains takes seven to eight
hours; for comparison, the per-domain models of \citet{gabar} take roughly two
hours each, so eight of them cost about sixteen hours and produce eight
artifacts rather than one.

The implementation uses PyTorch and PyTorch Geometric on Linux; a pinned
dependency specification ships with the code release.

\subsection{Renaming Invariance}
\label{app:invariance}

The claim that \gadar{} is invariant to renaming is enforced by construction
rather than learned, and it is checked by a unit test that runs on every
change to the representation: given a domain and a bijective renaming of its
predicates, action schemas, types, and objects, the test builds the graph for
the original and the renamed task and asserts that the node feature matrix,
the edge feature matrix, and the edge lists are identical. The test covers
both typing conventions in the benchmark suite, domains that declare PDDL
types and domains that spell types as ordinary unary predicates, and asserts
that the two produce the same structure after type compilation.

\ifappendixStandalone
\bibliography{aaai2027}
\end{document}
\fi
 just before \end{document} of gadar.tex
% The guard below detects which, so the same file serves both.
%
% Everything here is promised by an answer in the reproducibility checklist;
% the mapping is in Table~\ref{tab:checklist-map}.

\makeatletter
\@ifundefined{gadar}{
  \newif\ifappendixStandalone
  \appendixStandalonetrue
}{
  \newif\ifappendixStandalone
  \appendixStandalonefalse
}
\makeatother

\ifappendixStandalone
\documentclass[letterpaper]{article}
\usepackage{aaai2027}
\usepackage[hyphens]{url}
\usepackage{graphicx}
\urlstyle{rm}
\def\UrlFont{\rm}
\usepackage{natbib}
\usepackage{caption}
\usepackage{booktabs}
\usepackage{multirow}
\usepackage{amsfonts}
\usepackage{amsmath}
\usepackage{algorithm}
\usepackage{algorithmic}
\frenchspacing
\setlength{\pdfpagewidth}{8.5in}
\setlength{\pdfpageheight}{11in}
\setcounter{secnumdepth}{0}

\newcommand{\gadar}{GADAR}
\newcommand{\gadarbind}{\gadar{}-BIND}
\newcommand{\union}{UNION}

\title{GADAR: Supplementary Material}
\author{Anonymous Submission}
\affiliations{Anonymized for review}

\begin{document}
\maketitle
\fi

\appendix
\section{Supplementary Material}

This appendix supplies the detail promised by our reproducibility checklist
answers. Table~\ref{tab:checklist-map} maps each answer to the section that
discharges it.

\begin{table}[h]
\small
\centering
\caption{Checklist answers and where they are discharged.}
\label{tab:checklist-map}
\begin{tabular}{@{}llp{2.9cm}@{}}
\toprule
Item & Answer & Section \\
\midrule
1.1 outline / pseudocode        & yes     & \ref{app:pseudocode} \\
4.2 hyperparameter search       & partial & \ref{app:hyper} \\
4.5 code released on publication& yes     & \ref{app:code} \\
4.6 implementation comments     & partial & \ref{app:code} \\
4.7 seeding                     & yes     & \ref{app:seeds} \\
4.8 computing infrastructure    & partial & \ref{app:compute} \\
4.9 evaluation metrics          & yes     & \ref{app:metrics} \\
4.10 number of runs             & yes     & \ref{app:seeds} \\
4.13 final hyperparameters      & yes     & \ref{app:hyper} \\
\bottomrule
\end{tabular}
\end{table}

\subsection{Code and Data Availability}
\label{app:code}

The implementation is a single Python package built on PyTorch and PyTorch
Geometric, using PDDLGym to parse domains and problems and Fast Downward
\citep{fast_downward} to generate training plans. The full source, the
experiment configuration files that reproduce every number in the paper, and
the analysis scripts will be released publicly under a permissive license on
publication. They are not attached to this submission.

Each experiment is a YAML configuration file naming the training domains, the
featurization rung, and the evaluation domains; a run is a single invocation
of the entry point with one such file, so the configuration files are a
complete record of what was run. Files implementing the components introduced
here, in particular graph construction, the lifted domain layer, the binding
layer, and the chain layer, carry comments naming the paper section each step
corresponds to. Files inherited unchanged from the per-domain policy we build
on \citep{gabar} are not annotated in that way.

\subsection{Domains and Data}
\label{app:data}

We use the eight benchmark domains of \citet{gabar}, generated with
\citet{pddl-generators}. Training data is the set of states along a
satisficing plan for each small instance, labelled with the action the planner
took. Table~\ref{tab:data} gives the sizes. Instances the planner fails to
solve within its limit are skipped, so the collected count can be below the
number of instances the generator produced.

\begin{table}[h]
\small
\centering
\caption{Data per domain. Problems and states are training instances and the
labelled states collected from their plans; test problems are the domain's
full test set (our hard split).}
\label{tab:data}
\begin{tabular}{@{}lrrr@{}}
\toprule
Domain & Problems & States & Test problems \\
\midrule
Blocks     & 200 & 4{,}216 & 200 \\
Gripper    & 147 & 7{,}287 & 173 \\
Miconic    & 228 & 5{,}453 & 119 \\
Visitall   & 125 & 3{,}654 & 50 \\
Grid       & 192 & 6{,}880 & 48 \\
Logistics  & 156 & 6{,}732 & 96 \\
Spanner    & 234 & 3{,}801 & 96 \\
Rovers     & 312 & 3{,}001 & 54 \\
\midrule
Total      & 1{,}594 & 41{,}024 & 836 \\
\bottomrule
\end{tabular}
\end{table}

Ten percent of the states of each domain are held out as a validation split,
stratified by domain so every training domain is represented. The split is
deterministic given the seed. In the leave-one-domain-out experiments the
held-out domain contributes no training or validation data at all: its plans
are never collected for that run.

\subsection{Graph Construction}
\label{app:pseudocode}

Algorithm~\ref{alg:graph} states the construction of
$\mathcal{G} = (V, E, X, R)$ for one task, in the notation of the
representation section. Lines marked $\dagger$ are computed once per domain
and reused for every state of that domain; the rest are per state. The three
rungs of the ladder are prefixes of this procedure: \union{} runs lines
1 to 3 with symbol-identity features in place of $X$, \gadarbind{} runs lines
1 to 7, and \gadar{} runs all of it.

\begin{algorithm}[h]
\caption{Construction of the \gadar{} graph for $(s, G, D)$}
\label{alg:graph}
\begin{algorithmic}[1]
\STATE $V \leftarrow O \cup L \cup A \cup \{g\}$ where $L$ holds the ground
  literals of $s$ and of $G$, the latter flagged as goal atoms
\STATE add $E_{\mathrm{lit}}$: each literal to its argument objects, argument
  position on the edge
\STATE add $E_{\mathrm{act}}$: each schema to the objects of its
  \emph{applicable} groundings, parameter slot on the edge
\STATE $\dagger$ compile each declared PDDL type of $D$ into a unary static
  predicate, so typed and predicate-typed domains present one structure
\STATE $\dagger$ $V \leftarrow V \cup \Sigma \cup Q$: a node per predicate
  symbol of $D$, and an occurrence node $q_{a,k}$ per literal occurrence in
  the preconditions and effects of each schema $a$
\STATE $\dagger$ add $E_{\mathrm{occ}}$: each occurrence to its predicate
  symbol and its schema, carrying role $r(q)$ and argument-to-slot pairs
\STATE add $E_{\mathrm{inst}}$ and $E_{\mathrm{bind}}$: each literal to its
  predicate symbol, each object to its type symbol, and $(o, q_{a,k})$
  whenever an applicable grounding of $a$ puts $o$ at a slot that occurrence
  $k$ constrains
\STATE $\dagger$ add $E_{\mathrm{chain}}$: \textit{enables} between an add and
  a pre occurrence of a shared predicate, \textit{threatens} likewise for del
\STATE add $E_{\mathrm{eff}}$: from each applicable grounding, link its
  objects to any unsatisfied goal atom it adds and any true atom it deletes
\STATE compute $\mathrm{gdist}(a)$ by breadth-first search over
  \textit{enables} from the schemas adding a goal predicate; mark objects
  occurring in an open goal atom
\STATE emit $X$ and $R$ as indicator vectors; widths depend only on the global
  arity caps $(k_p, k_a)$, never on $|A|$, $|P|$, $|T|$ or $|O|$
\end{algorithmic}
\end{algorithm}

Decoding is unchanged from the per-domain policy: a GRU scores the schema
nodes against its hidden state, takes the best, then scores the object nodes
once per argument of the selected schema, feeding each choice back as the next
input. Search over these choices is a beam of width 2 over schemas and 3 over
objects per slot.

\subsection{Hyperparameters}
\label{app:hyper}

Table~\ref{tab:hyper} lists every setting used. Encoder and decoder
hyperparameters are inherited from \citet{gabar} without modification, which
is deliberate: the paper's claims are about the representation, and re-tuning
the network per rung would confound them. We therefore ran no hyperparameter
search over the architecture, and item 4.2 of the checklist is answered
``partial'' for that reason. The two settings we did choose are the global
arity caps $(k_p, k_a)$, set to the maxima over the training domains so that a
held-out domain with wider symbols still featurizes at full width, and the
training length, discussed below.

\begin{table}[h]
\small
\centering
\caption{Final hyperparameters. Identical across all three rungs except where
noted.}
\label{tab:hyper}
\begin{tabular}{@{}lr@{}}
\toprule
Setting & Value \\
\midrule
\multicolumn{2}{@{}l}{\emph{Encoder}} \\
Message-passing rounds $L$        & 9 \\
Representation size               & 64 \\
Attention heads                   & 1 \\
MLP layers per update block       & 2 \\
Global node                       & enabled \\
\midrule
\multicolumn{2}{@{}l}{\emph{Decoder}} \\
GRU layers                        & 3 \\
Beam width, schemas               & 2 \\
Beam width, objects per slot      & 3 \\
\midrule
\multicolumn{2}{@{}l}{\emph{Optimization}} \\
Optimizer                         & Adam \\
Learning rate                     & $5 \times 10^{-4}$ \\
Batch size                        & 64 \\
Dropout                           & 0.0 \\
Attention dropout                 & 0.0 \\
Weight decay                      & 0.0 \\
Epochs, structural rungs          & 500 \\
Epochs, \union{}                  & 100 \\
\midrule
\multicolumn{2}{@{}l}{\emph{Evaluation}} \\
Execution bound (steps)           & 500 \\
Search at test time               & none (greedy) \\
Revisit monitor                   & enabled \\
Training plan generator           & \texttt{fd-lama-first} \\
Plan-quality reference            & \texttt{fd-lama-first} \\
\bottomrule
\end{tabular}
\end{table}

\subsection{Seeds and Number of Runs}
\label{app:seeds}

A run is seeded once, and that seed initializes the network, the training
shuffle, the train and validation split, and the revisit monitor's fallback
sampling. The seed is recorded alongside every checkpoint, and the checkpoint
directory is keyed by the training domain set, the seed, the featurization,
and the resulting input widths, so runs that differ in any of these cannot
overwrite one another.

Each reported configuration is run for three seeds and the results averaged.
We report means without dispersion: three runs are too few for a meaningful
interval, and the differences the paper turns on, for instance 4\% against
68\% zero-shot coverage on the hard split of held-out Visitall, are far larger
than the spread we observe. No statistical test is applied, which is item 4.12
of the checklist.

\subsection{Evaluation Metrics}
\label{app:metrics}

Let $\mathcal{T}$ be a test set and $\mathrm{solved}(\pi) \subseteq
\mathcal{T}$ the instances on which policy $\pi$ reaches a goal state within
the execution bound. \emph{Coverage} is
\begin{equation*}
C(\pi) \;=\; \frac{|\mathrm{solved}(\pi)|}{|\mathcal{T}|} \times 100 .
\end{equation*}
Execution is greedy and search-free: at each state the highest-ranked
applicable action the decoder constructs is executed, and a rollout ends at
the goal, at a dead end, or at the bound. A rollout that proposes an
inapplicable action falls back to the next candidate in the beam; if none is
applicable the rollout fails.

\emph{Plan quality} compares plan lengths on the instances both the policy and
the reference planner solve. With $|\pi(t)|$ the policy's plan length on
instance $t$ and $|\rho(t)|$ the reference planner's,
\begin{equation*}
P(\pi) \;=\; \frac{\sum_{t \in \mathrm{solved}(\pi) \cap \mathrm{solved}(\rho)} |\rho(t)|}
                 {\sum_{t \in \mathrm{solved}(\pi) \cap \mathrm{solved}(\rho)} |\pi(t)|} .
\end{equation*}
Higher is better and $P > 1$ means plans shorter than the reference. $P$ is
conditioned on what a system solves, so a system with higher coverage is
scored on harder instances and $P$ must be read within a column rather than
across columns. Coverage is the primary metric for that reason.

Coverage is measured with a revisit monitor, which resamples among the
decoder's applicable proposals when a greedy rollout re-enters a state it has
already visited. The monitor is part of the execution harness and is applied
identically to every system, including the untrained control, which is what
makes the control informative: it measures what the harness achieves when the
learned component contributes nothing.

\subsection{Compute Environment}
\label{app:compute}

All training and all evaluation ran on a single NVIDIA A40 GPU per run, one
run per GPU. Training a single model on all eight domains takes seven to eight
hours; for comparison, the per-domain models of \citet{gabar} take roughly two
hours each, so eight of them cost about sixteen hours and produce eight
artifacts rather than one.

The implementation uses PyTorch and PyTorch Geometric on Linux; a pinned
dependency specification ships with the code release.

\subsection{Renaming Invariance}
\label{app:invariance}

The claim that \gadar{} is invariant to renaming is enforced by construction
rather than learned, and it is checked by a unit test that runs on every
change to the representation: given a domain and a bijective renaming of its
predicates, action schemas, types, and objects, the test builds the graph for
the original and the renamed task and asserts that the node feature matrix,
the edge feature matrix, and the edge lists are identical. The test covers
both typing conventions in the benchmark suite, domains that declare PDDL
types and domains that spell types as ordinary unary predicates, and asserts
that the two produce the same structure after type compilation.

\ifappendixStandalone
\bibliography{aaai2027}
\end{document}
\fi
